% Source: https://artofproblemsolving.com/community/c3h1964721_sample_packet_1
% Stable identity: aops-topic-1964721-q01
An affine map (a linear map followed by a translation) carries any triangle to any other, preserves midpoints, centroids and the ratio of areas of figures. So we may compute with the convenient triangle $A=(0,0)$, $B=(6,0)$, $C=(0,6)$, of area $18$. Then $O=(2,2)$, $D=(3,3)$, $E=(3,0)$, $F=(0,3)$, and the six centroids are
\[
\begin{aligned}
&AFO:\ \left(\dfrac{2}{3},\dfrac{5}{3}\right),\quad FOC:\ \left(\dfrac{2}{3},\dfrac{11}{3}\right),\quad COD:\ \left(\dfrac{5}{3},\dfrac{11}{3}\right),\\
&DOB:\ \left(\dfrac{11}{3},\dfrac{5}{3}\right),\quad BOE:\ \left(\dfrac{11}{3},\dfrac{2}{3}\right),\quad EOA:\ \left(\dfrac{5}{3},\dfrac{2}{3}\right).
\end{aligned}
\]
Multiplying all coordinates by $3$ (which multiplies areas by $9$) gives the convex hexagon $(2,5),(2,11),(5,11),(11,5),(11,2),(5,2)$. By the shoelace formula its area is
\[
\dfrac{1}{2}\left|12-33-96-33+12+21\right|=\dfrac{117}{2},
\]
so the hexagon itself has area $\dfrac{117}{18}=\dfrac{13}{2}$. The ratio is
\[
\dfrac{\dfrac{13}{2}}{18}=\boxed{\dfrac{13}{36}} .
\]

\begin{asy}
usepackage("amsmath");
size(6cm);
defaultpen(fontsize(10pt));
pair A=(1.2000,4.5000),B=(0.0000,0.0000),C=(6.0000,0.0000),O=(2.4000,1.5000),D=(3.0000,0.0000),Ep=(0.6000,2.2500),Fp=(3.6000,2.2500);
fill((2.4000,2.7500)--(4.0000,1.2500)--(3.8000,0.5000)--(1.8000,0.5000)--(1.0000,1.2500)--(1.4000,2.7500)--cycle,palegray);draw((2.4000,2.7500)--(4.0000,1.2500)--(3.8000,0.5000)--(1.8000,0.5000)--(1.0000,1.2500)--(1.4000,2.7500)--cycle,blue+linewidth(0.9));
dot((2.4000,2.7500),blue+2.5);dot((4.0000,1.2500),blue+2.5);dot((3.8000,0.5000),blue+2.5);dot((1.8000,0.5000),blue+2.5);dot((1.0000,1.2500),blue+2.5);dot((1.4000,2.7500),blue+2.5);
draw(A--B--C--cycle,linewidth(1));draw(A--D);draw(C--Ep);draw(B--Fp);
dot(A);dot(B);dot(C);dot(D);dot(Ep);dot(Fp);dot(O);
label("$A$",A,N);label("$B$",B,SW);label("$C$",C,SE);label("$D$",D,S);label("$E$",Ep,W);label("$F$",Fp,NE);label("$O$",O,0.6*S+0.9*E);
label("\mbox{}",(min(currentpicture,user=true)+max(currentpicture,user=true))/2,black); // keeps the text after the figure black when LaTeX typesets the labels inline
\end{asy}
